\part{Engineering: building, debugging, improving}

\chapter{The build system}

\begin{goals}
  \item Read the project Makefile and understand every rule
  \item Understand why there are two compilers
  \item Know the weaknesses of this build and how to fix them
\end{goals}

\section{Two compilers, one project}

\filetag{lytlnybl/Makefile}
\begin{lstlisting}[style=mk]
CC = i686-elf-gcc          # for code that runs on lytlnyblOS
HOST_CC = gcc              # for mkfs, which runs on your machine
LD = ld
NASM = nasm
OBJCOPY = objcopy
AR = i686-elf-ar

CFLAGS = -g -m32 -Iinclude -ffreestanding
USER_CFLAGS = $(CFLAGS) -fno-pie -fno-pic
\end{lstlisting}

\cod{-fno-pie -fno-pic} for user programs is worth explaining. Modern compilers default to
generating \emph{position-independent code}, which can be loaded anywhere --- useful for address
space randomisation on a real OS. But our loader maps every program at exactly \hexa{400000} and
there is no dynamic relocation, so position independence is pure overhead and can produce
constructs the loader would have to handle. Turning it off gives plain absolute code.

\section{Finding sources automatically}

\begin{lstlisting}[style=mk]
KERNEL_C_FILES = $(shell find kernel memory filesystem tasks -name '*.c')
KERNEL_ASM_FILES = $(shell find kernel memory tasks -name '*.asm')

KERNEL_C_OBJECTS = $(KERNEL_C_FILES:%.c=$(BUILD_DIR)/%.o)
KERNEL_ASM_OBJECTS = $(KERNEL_ASM_FILES:%.asm=$(BUILD_DIR)/%.asm.o)
\end{lstlisting}

\cod{\$(shell find ...)} runs a command at parse time and captures its output, so adding a new
\cod{.c} file to \filepath{kernel/} requires no Makefile change. The substitution
\cod{\$(FILES:\%.c=\$(BUILD\_DIR)/\%.o)} turns \cod{kernel/interrupts.c} into
\cod{build/kernel/interrupts.o}, mirroring the source tree inside \filepath{build/}.

\section{Pattern rules}

\begin{lstlisting}[style=mk]
$(BUILD_DIR)/%.o: %.c
	mkdir -p $(dir $@)
	$(CC) $(CFLAGS) -c $< -o $@

$(BUILD_DIR)/%.asm.o: %.asm
	mkdir -p $(dir $@)
	$(NASM) -f elf32 -g -F dwarf $< -o $@
\end{lstlisting}

\begin{center}
\begin{tabularx}{0.9\textwidth}{@{}l X@{}}
\toprule
\textbf{Variable} & \textbf{Means} \\
\midrule
\cod{\$@} & The target being built \\
\cod{\$<} & The first prerequisite \\
\cod{\$\^{}} & All prerequisites \\
\cod{\$(dir \$@)} & The directory part of the target \\
\bottomrule
\end{tabularx}
\end{center}

\cod{-f elf32} tells NASM to produce a relocatable object file rather than a flat binary, so the
linker can combine it with the C objects. \cod{-g -F dwarf} adds debug information in the format GDB
expects --- which is what makes stepping through assembly in the debugger possible.

\section{The critical link order}

\begin{lstlisting}[style=mk]
$(KERNEL_ELF): $(KERNEL_C_OBJECTS) $(KERNEL_ASM_OBJECTS)
	$(LD) -m elf_i386 -T $(KERNEL_LINKER) \
		$(KERNEL_ENTRY_OBJECT) \
		$(KERNEL_C_OBJECTS) \
		$(filter-out $(KERNEL_ENTRY_OBJECT),$(KERNEL_ASM_OBJECTS)) \
		-o $@
\end{lstlisting}

\begin{concept}{Why \cod{kernel\_intro.asm.o} must be first}
The linker lays out \cod{.text} in the order the object files appear. The linker script puts
\cod{.text} at \hexa{9000}, and the bootloader jumps to \hexa{9000}.

Therefore \hi{whatever code is first in the link line ends up at the entry point}.
\cod{\$(KERNEL\_ENTRY\_OBJECT)} is listed explicitly first, and \cod{filter-out} removes it from the
rest so it is not linked twice.

If that ordering broke, the bootloader would jump into the middle of some random function, in 16-bit
mode, with no diagnostic whatsoever. Note the linker already warns about a related issue:
\begin{lstlisting}[style=txt]
ld: warning: cannot find entry symbol start; defaulting to 00009000
\end{lstlisting}
because \cod{start:} in the assembly file is not declared \cod{global}. It works only because of the
link order. Adding \cod{global start} would make the intent explicit and the build robust.
\end{concept}

\section{From ELF to a disk image}

\begin{lstlisting}[style=mk]
$(KERNEL_BIN): $(KERNEL_ELF)
	$(OBJCOPY) -O binary $< $@
\end{lstlisting}

An ELF file has headers, a symbol table, section metadata and debug information. The bootloader
understands none of that --- it copies raw sectors to \hexa{9000} and jumps. \cod{objcopy -O binary}
strips everything and leaves just the bytes, in the order the linker placed them.

Both files are kept: \cod{kernel.bin} goes into the image, and \cod{kernel.elf} stays behind for GDB,
which needs the symbols.

\section{What is missing from this build}

\begin{danger}{No header dependencies}
This is the most practically annoying weakness. The rule says \cod{\$(BUILD\_DIR)/\%.o: \%.c} ---
the object depends on the \cod{.c} file and nothing else.

So editing \filepath{include/memory/vmm.h} rebuilds \textbf{nothing}. You change a struct, run
\cod{make}, it says ``nothing to be done'', and you spend the next hour debugging a binary built
from two incompatible versions of the same header.

The fix is two lines:
\begin{lstlisting}[style=mk]
CFLAGS += -MMD -MP
-include $(shell find $(BUILD_DIR) -name '*.d' 2>/dev/null)
\end{lstlisting}
\cod{-MMD} makes GCC emit a \cod{.d} file listing every header the source included, and \cod{-include}
pulls those in as extra rules. After that, touching a header rebuilds exactly what depends on it.
\end{danger}

\begin{warn}{No warnings enabled}
\cod{CFLAGS} has no \cod{-Wall -Wextra}. Turning them on for this project produces only a handful of
findings --- a few unused variables and one intentional switch fall-through --- so the cost of
adopting them is almost zero and the future benefit is large.

Be aware of the limits, though: \cod{-Wall -Wextra} does \emph{not} catch the uninitialised
\cod{word\_count} in the shell, because taking its address defeats the analysis. Warnings are a
floor, not a ceiling.
\end{warn}

\begin{warn}{Host linker, and generated binaries in git}
\cod{LD = ld} uses the host linker with \cod{-m elf\_i386}. It works on an x86-64 Linux host but
would fail on, say, an ARM machine where the host \cod{ld} cannot target i386. \cod{i686-elf-ld}
would be correct.

Separately, \filepath{rootfs/bin/shell} and \filepath{rootfs/bin/user\_test} are tracked in
git even though they are build outputs. They show as modified after every build, which clutters every
diff. They belong in \cod{.gitignore}.
\end{warn}

% ============================================================
\chapter{Debugging a kernel}

\begin{goals}
  \item Know the four techniques available and when each applies
  \item Be able to use GDB against a running kernel
  \item Be able to read a QEMU interrupt trace
\end{goals}

\section{Four techniques}

\begin{center}
\begin{tabularx}{\textwidth}{@{}l l X@{}}
\toprule
\textbf{Technique} & \textbf{Cost} & \textbf{Best for} \\
\midrule
Print to screen & trivial & ``Does it reach this line?'' ``What is this value?'' \\
GDB via QEMU & medium & Inspecting state at a precise point; stepping \\
QEMU tracing & low & Faults, resets, interrupt storms --- things that happen too fast to see \\
Reading the code & free & Everything else, and you will do more of this than you expect \\
\bottomrule
\end{tabularx}
\end{center}

\section{Printing}

Crude and extremely effective. The kernel does it itself before any driver exists:

\begin{lstlisting}[style=c]
volatile char* vga = (volatile char*)0xB8000;
vga[0] = 'C'; vga[1] = 0x02;
\end{lstlisting}

Two bytes and you know execution reached that line. Using different screen positions for different
checkpoints gives a crude progress display that survives even a completely broken driver.

\begin{warn}{But printing can be the thing that is broken}
Chapter 20 showed the failure mode: if \cod{terminal} is corrupt, every print attempt faults. When
the system dies \emph{during} your debug output, printing tells you nothing. That is exactly when
the next two techniques earn their keep.
\end{warn}

\section{GDB against QEMU}

QEMU can expose a GDB server and start paused:

\filetag{lytlnybl/Makefile}
\begin{lstlisting}[style=mk]
run: $(KERNEL_IMG)
	qemu-system-i386 -drive format=raw,file=$(KERNEL_IMG) #-s -S
\end{lstlisting}

Uncomment \cod{-s -S}: \cod{-s} opens a GDB server on port 1234, \cod{-S} freezes the CPU before the
first instruction. Then, in another terminal:

\begin{lstlisting}[style=sh]
gdb build/kernel.elf
(gdb) target remote localhost:1234
(gdb) break kernel_main
(gdb) continue
\end{lstlisting}

The project ships a \filepath{.gdbinit} to automate the connection.

\subsection{Commands worth knowing}

\begin{center}
\begin{tabularx}{\textwidth}{@{}l X@{}}
\toprule
\textbf{Command} & \textbf{Does} \\
\midrule
\cod{info registers} & Show every register, including \reg{CR0}, \reg{CR2}, \reg{CR3} \\
\cod{x/16xb 0xB8000} & Dump 16 bytes in hex from video memory \\
\cod{x/10i \$eip} & Disassemble the next 10 instructions \\
\cod{p/x terminal} & Print a struct in hex \\
\cod{watch terminal.width} & \textbf{Break when a variable changes} \\
\cod{bt} & Backtrace --- the chain of calls \\
\cod{stepi} / \cod{nexti} & Step one machine instruction \\
\bottomrule
\end{tabularx}
\end{center}

\begin{concept}{\cod{watch} is the one that finds memory corruption}
When a variable is being overwritten by something and you have no idea what, a hardware watchpoint
is the tool. \cod{watch terminal.width} tells the processor to trap on any write to that address,
and GDB stops you at the exact instruction responsible --- which may be in a completely unrelated
file.

For the bug in the next chapter, \cod{watch terminal.width} would have pointed straight at a
\cod{memcpy} inside the program loader.
\end{concept}

\section{QEMU tracing}

For faults that happen faster than you can observe, QEMU can log every interrupt and every CPU
reset:

\begin{lstlisting}[style=sh]
qemu-system-i386 -drive format=raw,file=build/kernel.img \
    -display none -no-reboot \
    -d int,cpu_reset -D trace.log
\end{lstlisting}

\begin{center}
\begin{tabularx}{0.95\textwidth}{@{}l X@{}}
\toprule
\textbf{Option} & \textbf{Effect} \\
\midrule
\cod{-d int} & Log every interrupt and exception with full register state \\
\cod{-d cpu\_reset} & Log every processor reset \\
\cod{-no-reboot} & \textbf{Stop instead of rebooting} on triple fault \\
\cod{-D file} & Write the log to a file \\
\bottomrule
\end{tabularx}
\end{center}

\cod{-no-reboot} is the important one. Without it, a triple fault silently restarts the machine and
you see a boot loop with no information. With it, QEMU prints \cod{Triple fault} and exits, leaving
the log ending exactly at the failure.

A log line looks like this:

\begin{lstlisting}[style=txt]
   278: v=0e e=0002 i=0 cpl=0 IP=0008:0000ad98 pc=0000ad98 SP=0010:c0000000 CR2=bffffffc
\end{lstlisting}

\begin{center}
\begin{tabularx}{0.95\textwidth}{@{}l X@{}}
\toprule
\textbf{Field} & \textbf{Meaning} \\
\midrule
\cod{278} & Sequence number --- how many interrupts have occurred \\
\cod{v=0e} & Vector \hexa{0E} = 14 = page fault \\
\cod{e=0002} & Error code: bit 1 set = it was a write, bit 0 clear = page not present \\
\cod{cpl=0} & Current privilege level: the kernel \\
\cod{IP=0008:0000ad98} & Where it faulted \\
\cod{SP=0010:c0000000} & Stack pointer at the time \\
\cod{CR2=bffffffc} & The address that was accessed \\
\bottomrule
\end{tabularx}
\end{center}

\section{Reading the screen without a screen}

One more trick, useful for automated testing. QEMU's monitor can capture the framebuffer:

\begin{lstlisting}[style=sh]
qemu-system-i386 -drive format=raw,file=build/kernel.img \
    -display none -monitor unix:/tmp/mon.sock,server,nowait
\end{lstlisting}

Connect to that socket and issue \cod{screendump /tmp/shot.ppm}. You get a 720$\times$400 image of
the text screen.

\begin{warn}{\cod{pmemsave} will not read video memory}
It is tempting to dump \hexa{B8000} directly with the monitor's \cod{pmemsave} command. It returns
zeros: that address is a device region, not RAM, and the command only reads RAM.

\cod{screendump} works because it asks the display device for its rendered output. Since the VGA
font is a known 8$\times$16 bitmap (extractable from the \cod{vgabios.bin} that ships with QEMU),
each 9$\times$16 cell of the image can be matched back to a character --- giving you the text of the
screen as a string, from a script, with no display attached. That is how the boot output quoted
throughout this course was captured.
\end{warn}

% ============================================================
\chapter{Case study: a reboot loop}

\begin{goals}
  \item Follow a complete diagnosis from symptom to root cause
  \item See how to work backwards along a chain of causes
  \item Learn the method, not just this answer
\end{goals}

This chapter walks through a real fault in this project, diagnosed on a real machine. The
intermediate observations are exactly those produced by the tools, not reconstructed afterwards.

\section{The symptom}

A clean rebuild of the project boots, prints its self-test output, reaches the line
\cod{Running shell:} --- and then the screen goes blank and the BIOS banner appears again. And
again. A reboot loop, with no error message.

Curiously, the \emph{prebuilt} image already in the repository boots fine and reaches a working
shell prompt. Same source, same commands, different behaviour.

\begin{concept}{First observation: ``it works for me'' is data}
Two builds of identical source behaving differently is not noise. It says the fault depends on
something that varies between builds --- almost always \emph{layout}: addresses, sizes, alignment.
That single observation already eliminates whole categories of bug (logic errors, race conditions)
and points at memory.
\end{concept}

\section{Step 1: catch the fault instead of the reboot}

\begin{lstlisting}[style=sh]
qemu-system-i386 -drive format=raw,file=build/kernel.img \
    -display none -no-reboot -d int,cpu_reset -D trace.log
\end{lstlisting}

QEMU now exits with \cod{Triple fault} instead of rebooting. The tail of the log:

\begin{lstlisting}[style=txt]
   278: v=0e e=0002 i=0 cpl=0 IP=0008:0000ad98 SP=0010:c0000000 CR2=bffffffc
   ...
   279: v=08 e=0000 i=0 cpl=0 IP=0008:0000ad98 SP=0010:c0000000
   ...
check_exception old: 0x8 new 0xe
Triple fault
\end{lstlisting}

Reading it with the table from the previous chapter:

\begin{itemize}
  \item \cod{v=0e}: a page fault, \cod{e=0002} means a write to a non-present page.
  \item \cod{CR2=bffffffc}: the address accessed was \hexa{BFFFFFFC}.
  \item \cod{SP=c0000000}: the stack pointer was \hexa{C0000000}.
  \item Then \cod{v=08}: a double fault. Then triple fault.
\end{itemize}

\begin{concept}{Step 1 conclusion}
\reg{ESP} is \hexa{C0000000} and the faulting address is \hexa{BFFFFFFC} --- exactly four bytes
below. A push. \hi{This is a stack overflow}, matching the signature from Chapter 28
precisely.

And \hexa{C0000000} is \cod{HEAP\_START}. So a stack allocated from the kernel heap has grown down
past the very beginning of the heap.
\end{concept}

\section{Step 2: find out where}

\begin{lstlisting}[style=sh]
i686-elf-nm -n build/kernel.elf | grep " [tT] "
\end{lstlisting}

Matching \cod{IP=0000ad98} against the symbol table gives \cod{vga\_text\_writeline + 0x0} --- the
function's very first instruction, \cod{push ebp}. Consistent: the stack was already exhausted, and
the first push of the next call tipped it over.

But that only says \emph{where it ran out}, not \emph{why}. For that, look further back in the log.

\section{Step 3: the pattern before the crash}

\begin{lstlisting}[style=txt]
   240: v=00 e=0000 i=0 cpl=0 IP=0008:0000ad47 SP=0010:c00016ec
   241: v=00 e=0000 i=0 cpl=0 IP=0008:0000ad47 SP=0010:c0001650
   242: v=00 e=0000 i=0 cpl=0 IP=0008:0000ad47 SP=0010:c00015b4
   ...
   277: v=00 e=0000 i=0 cpl=0 IP=0008:0000ad47 SP=0010:c0000060
   278: v=0e ...                                SP=0010:c0000000
\end{lstlisting}

This is the answer, and it is unmistakable:

\begin{itemize}
  \item \cod{v=00} is \hi{exception 0: divide by zero}.
  \item The same instruction address, \hexa{AD47}, every time.
  \item \reg{ESP} decreasing by exactly \hexa{9C} (156 bytes) each round.
  \item Exactly 107 repetitions (log entries 171 to 277) before the stack is exhausted ---
        and $107 \times \hexa{9C} = 16{,}692$ bytes, which is the 16 KB kernel stack.
\end{itemize}

\begin{concept}{Step 3 conclusion}
An infinite recursion. Something divides by zero; the handler tries to print the error; printing
divides by zero again. Each round consumes 156 bytes of stack and never returns. Exactly the
scenario described in Chapter 20 --- now observed rather than predicted.
\end{concept}

\section{Step 4: what is being divided}

\begin{lstlisting}[style=sh]
i686-elf-objdump -d -S build/kernel.elf --start-address=0xad30 --stop-address=0xad80
\end{lstlisting}

\begin{lstlisting}[style=txt]
    ad39:  mov 0x8(%ebp),%eax
    ad3c:  mov 0x8(%eax),%ecx      ; ECX = terminal->width
    ad3f:  mov -0xc(%ebp),%eax     ; EAX = pos
    ad42:  mov $0x0,%edx
    ad47:  div %ecx                ; <-- the faulting instruction
\end{lstlisting}

Which is this source line:

\begin{lstlisting}[style=c]
terminal->row = pos / terminal->width;
\end{lstlisting}

So \cod{terminal->width} is zero. But \cod{vga\_text\_init} sets it to 80 during boot, and the
system clearly printed many lines successfully before failing. \hi{Therefore something
overwrote it.}

\section{Step 5: who could overwrite it}

\begin{lstlisting}[style=sh]
i686-elf-nm build/kernel.elf | grep -E " (B|b) |_kernel_end"
\end{lstlisting}

\begin{lstlisting}[style=txt]
00013188 B terminal
00013400 B tss
00013468 B _kernel_end
\end{lstlisting}

Now the arithmetic from Chapter 25:

$$\hexa{13468} \div \hexa{1000} = 19.27 \rightarrow 19$$

The protection loop is \cod{for (i = 9; i < 19; i++)}, protecting frames 9--18, that is
\hexa{9000}--\hexa{12FFF}. \hi{Frame 19 --- \hexa{13000}--\hexa{13FFF} --- is left free.} And
\cod{terminal} is at \hexa{13188}, inside it.

The chain is now complete:

\begin{center}
\begin{tikzpicture}[font=\sffamily\scriptsize,
  n/.style={draw=grayborder,rounded corners=2pt,minimum width=10.6cm,minimum height=0.62cm,align=left,inner xsep=8pt}]
  \node[n,fill=accentlight,draw=accent] (a) {1. \texttt{\_kernel\_end} is not page-aligned; the division truncates};
  \node[n,fill=accentlight,draw=accent,below=0.2cm of a] (b) {2. The last partial page of the kernel is not marked used};
  \node[n,fill=amberlight,draw=amber,below=0.2cm of b] (c) {3. \texttt{alloc\_frame()} hands out frame 19 to the program loader};
  \node[n,fill=amberlight,draw=amber,below=0.2cm of c] (d) {4. \texttt{memcpy} writes the shell binary over the kernel's globals};
  \node[n,fill=redlight,draw=red,below=0.2cm of d] (e) {5. \texttt{terminal.width} becomes 0};
  \node[n,fill=redlight,draw=red,below=0.2cm of e] (f) {6. The next print divides by zero};
  \node[n,fill=redlight,draw=red,below=0.2cm of f] (g) {7. The handler prints, which divides by zero, forever};
  \node[n,fill=redlight,draw=red,below=0.2cm of g] (h) {8. Kernel stack exhausted $\rightarrow$ page fault $\rightarrow$ double $\rightarrow$ triple fault $\rightarrow$ reboot};
  \foreach \x/\y in {a/b,b/c,c/d,d/e,e/f,f/g,g/h} \draw[-{Stealth[length=2mm]},gray] (\x)--(\y);
\end{tikzpicture}
\end{center}

\section{Step 6: confirm by fixing}

A hypothesis is not a diagnosis until the fix proves it. One line:

\begin{lstlisting}[style=c]
uint32_t kernel_frame_end = ((uint32_t)&_kernel_end + FRAME_SIZE - 1) / FRAME_SIZE;
\end{lstlisting}

Rebuild, boot:

\begin{lstlisting}[style=txt]
Welcome to the lytlnyblOS!
Kernel multiprocessing tests:
NEW KERNEL PROCESS RUNNING
BACK IN MAIN PROCESS
Memory allocation tests:
Values: 10 20 30 40 50
Filesystem tests:
Root directory contents:
bin file.txt notes.txt
Bin directory contents:
user_test shell
Running shell:
>
\end{lstlisting}

Stable, with a working prompt, and still stable ten seconds later. Confirmed.

\section{Why the prebuilt image was fine}

Because \cod{\_kernel\_end} lands somewhere different in every build. In the repository's prebuilt
image, \cod{kernel.bin} is 22,920 bytes; in the clean rebuild it is 22,888. A 32-byte difference ---
from nothing more than toolchain version --- moved which globals fall in the unprotected last page.

\begin{concept}{The general lesson}
\hi{The bug was always there.} It was present in every build, including the ones that worked
perfectly. What varied was only whether it had visible consequences.

This is the defining property of latent memory bugs, and why ``it works'' is such weak evidence in
systems programming. The fix is not to test more; it is to remove the class of error --- here, by
always rounding up when converting a byte address to a page index.
\end{concept}

\begin{recipe}{reproducing the reboot loop}{ch45-reboot-loop}
Injects the fault deliberately --- a kernel process sets \cod{terminal.width = 0} and then prints
--- so the failure is reproducible instead of depending on where \cod{\_kernel\_end} happens to
land.

\cod{make trace} is the target to use. It boots with \cod{-no-reboot -d int,cpu\_reset}, stops at
the triple fault instead of restarting, and prints the last exceptions. You get the signature to
recognise: a run of \cod{v=00} at one fixed instruction address with \reg{SP} falling by exactly
\hexa{9C} each time.

The addresses will not match the ones quoted above --- the fault is injected at a different point
in boot, on a different stack --- but the recursion signature is identical, and that is the part
worth being able to spot.
\end{recipe}

\section{The method, generalised}

\begin{enumerate}
  \item \textbf{Make the failure stop instead of hiding.} \cod{-no-reboot} turns an
        information-free reboot loop into a log that ends at the fault.
  \item \textbf{Read the last fault precisely.} Vector, error code, \reg{CR2}, \reg{ESP}. Each is a
        strong constraint.
  \item \textbf{Look for a pattern before it.} A repeating vector at a fixed address with a
        monotonically moving \reg{ESP} is recursion. One-off faults are different bugs entirely.
  \item \textbf{Map addresses back to source.} \cod{nm} and \cod{objdump} turn a hex address into a
        line of C.
  \item \textbf{Ask what could have written that value} rather than what read it. The reader is
        rarely the culprit.
  \item \textbf{Prove it with a fix.} If the minimal fix does not eliminate the symptom, the
        diagnosis is wrong, however convincing it sounded.
\end{enumerate}

% ============================================================
\chapter{Known issues and how to fix them}

\begin{goals}
  \item Have a single list of every weakness identified in this course
  \item Know the relative priority of each
  \item Have a concrete starting point for contributing
\end{goals}

Everything below was identified while writing this course, by reading the code and by running the
system. They are grouped by severity. None of them makes the project less valuable as teaching
material --- several of them make it \emph{more} valuable, because working through the fix teaches
more than reading correct code would.

\section{Serious: correctness and safety}

\begin{center}
\begin{tabularx}{\textwidth}{@{}p{3.3cm} X p{2.9cm}@{}}
\toprule
\textbf{Issue} & \textbf{Effect} & \textbf{Where} \\
\midrule
Unvalidated user pointers & A user program can make the kernel write a zero byte to any address, including kernel memory & \cod{interrupts.c}, \cod{SYSCALL\_WRITE} \\[3pt]
Write one byte past \cod{count} & Corrupts the caller's buffer; triggered by the shell's \cod{cat} & \cod{interrupts.c} \\[3pt]
\cod{load\_program} return not checked & \cod{run} on a missing file hangs the shell forever & \cod{interrupts.c}, \cod{FS\_RUN} \\[3pt]
Page index rounding & Latent: the kernel's last partial page could be handed out & \cod{pmm.c} --- \textbf{fixed}, see Ch.\ 45 \\[3pt]
VLA on the kernel stack & Programs over $\sim$14 KB overflow the 16 KB kernel stack & \cod{loader.c} \\[3pt]
\cod{destroy\_process} head removal & A terminated head process is never unlinked & \cod{procman.c} \\[3pt]
\cod{unmap\_page} on the kernel directory & Punches a permanent hole in the identity map shared by all processes & \cod{procman.c} \\[3pt]
No \cod{free\_frame} on teardown & Physical memory leaks on every process exit & \cod{procman.c} \\[3pt]
\cod{kfree} of the running kernel stack & Use-after-free during \cod{exit} & \cod{context.c} \\
\bottomrule
\end{tabularx}
\end{center}

\section{Moderate: robustness}

\begin{center}
\begin{tabularx}{\textwidth}{@{}p{3.3cm} X p{2.9cm}@{}}
\toprule
\textbf{Issue} & \textbf{Effect} & \textbf{Where} \\
\midrule
No idle process & \cod{get\_next\_process} can return a non-runnable process; the kernel busy-waits instead of halting & \cod{scheduler.c} \\[3pt]
\cod{\_start} vs \cod{start} & The shell's entry point works only by link order & \cod{shell.c}, \cod{user.ld} \\[3pt]
Uninitialised \cod{word\_count} & Shell behaviour depends on stack garbage & \cod{shell.c} \\[3pt]
\cod{tokenize\_line} capacity & More than 8 words overflows a heap allocation & \cod{shell.c} \\[3pt]
Word strings never freed & The shell leaks memory per command & \cod{shell.c} \\[3pt]
\cod{read} result ignored & \cod{strlen} on a non-terminated buffer & \cod{shell.c} \\[3pt]
User stack is one page & Deep recursion kills the process instead of growing the stack & \cod{procman.c} \\[3pt]
Only 4 MB identity-mapped & A ceiling on total allocation; faults if crossed & \cod{vmm.c} \\[3pt]
No ATA timeouts & A stuck disk hangs the kernel forever & \cod{ata.c} \\[3pt]
Global descriptor table & Descriptors are shared between processes and leak on exit & \cod{fs\_manager.c} \\
\bottomrule
\end{tabularx}
\end{center}

\section{Minor: hygiene}

\begin{itemize}
  \item No \cod{-MMD} header dependencies in the Makefile --- editing a header rebuilds nothing.
  \item No \cod{-Wall -Wextra}.
  \item \cod{LD = ld} rather than \cod{i686-elf-ld}.
  \item \cod{rootfs/bin/*} tracked in git although generated.
  \item \cod{\_\_attribute\_\_((packed))} on \cod{process\_t}, where it is unnecessary and harmful.
  \item \cod{syscall\_entry} pushes the interrupt number and error code in the opposite order to the
        ISR stubs.
  \item \cod{timer\_wait\_ms} counts ticks, not milliseconds, despite its name.
  \item The intentional fall-through in \cod{keyboard\_handler} is not marked as such.
  \item \cod{vga\_text\_write\_dec(t, 0)} prints the \cod{'0'} but does not advance the cursor, so
        the next thing written covers it and the zero disappears. Every other value goes through a
        loop that advances correctly. Found by running the Chapter 27 example, and hit again in the
        Chapter 39 one.
  \item No scrolling: the screen wraps instead.
  \item VGA text wraps at 60 sectors' worth of kernel with no build-time check.
\end{itemize}

\section{A suggested order of work}

If you want to improve the project, this order gives the most benefit for the least risk:

\begin{enumerate}
  \item \textbf{\cod{-MMD} in the Makefile.} Five minutes, and it removes an entire class of
        confusing rebuild bugs. Do this before touching anything else.
  \item \textbf{The three-line NULL check in \cod{FS\_RUN}.} Turns a hang into an error message.
  \item \textbf{Round up in \cod{pmm.c}.} Already done --- see Chapter 45 --- but read the
        diagnosis, because the same truncation pattern can appear anywhere a byte address is
        converted to a page or block index.
  \item \textbf{A \cod{copy\_from\_user} / \cod{copy\_to\_user} pair.} The highest-value change in
        the codebase; closes the whole unvalidated-pointer category.
  \item \textbf{An idle process.} Makes the scheduler correct by construction and lets the CPU
        actually sleep.
  \item \textbf{Fix \cod{destroy\_process}} --- unlink correctly, stop unmapping from
        \cod{kernel\_directory}, free the frames, defer the \cod{kfree} of the running stack.
  \item \textbf{Stream the program loader} page by page instead of the VLA.
  \item \textbf{Shell hygiene}: initialise the counter, enforce capacity, free the words, use
        \cod{read}'s return value.
\end{enumerate}

% ============================================================
\chapter{Where to go from here}

\begin{goals}
  \item Know which extensions are worth attempting and in what order
  \item Understand what each one teaches
\end{goals}

\section{Extensions, by difficulty}

\subsection{Approachable}

\begin{description}
  \item[Screen scrolling] Move 24 rows up with \cod{memcpy} and clear the last. Teaches: the VGA
    buffer layout, and why terminals behave as they do.
  \item[Arrow keys] The extended-scancode path already exists and does nothing. Fill it in, then
    add a command history to the shell. Teaches: scancode handling and line editing.
  \item[More commands] \cod{help}, \cod{clear}, \cod{echo}. Requires first relaxing the
    \cod{word\_count != 2} restriction. Teaches: argument parsing.
  \item[Buffered \cod{printf}] Accumulate and flush on newline. Teaches: why real libraries buffer,
    and how much system calls cost.
\end{description}

\subsection{Moderate}

\begin{description}
  \item[Growing the user stack on demand] Catch a page fault just below the stack region and map a
    new page. Fifteen lines that turn a crash into a feature. Teaches: what page faults are really
    for.
  \item[Indirect blocks] An eleventh inode pointer that points at a block of pointers. Raises the
    file size limit from 5 KB to $\sim$69 KB. Teaches: how real filesystems scale.
  \item[\cod{.} and \cod{..}, and a working directory] Teaches: path resolution and per-process
    state.
  \item[Keyboard ring buffer] Move the interrupt handler's writes into a kernel buffer and copy to
    user space in process context. Teaches: why interrupt handlers should be short, and removes the
    \reg{CR3} manipulation.
  \item[Priorities in the scheduler] Give interactive processes shorter, more frequent slices.
    Teaches: what schedulers actually do.
\end{description}

\subsection{Ambitious}

\begin{description}
  \item[ELF loading] Parse real ELF headers instead of flat binaries. Teaches: how program loading
    works everywhere else, and enables separate read-only text and writable data.
  \item[Separate \cod{fork} and \cod{exec}] Requires copying an address space. Copy-on-write makes
    it efficient and is a beautiful use of page faults. Teaches: the core of Unix process creation.
  \item[Disk interrupts instead of polling] Frees the CPU during I/O. Teaches: asynchronous I/O and
    why it complicates everything.
  \item[Higher-half kernel] Move the kernel to \hexa{C0000000} and give user space the whole lower
    range. Teaches: how real systems lay out address spaces, and a great deal about paging.
  \item[Pipes] \cod{ls / | cat}. Requires a kernel buffer, blocking reads and writes, and shell
    syntax. Teaches: inter-process communication.
\end{description}

\section{A closing thought}

The system in this course is about 5,500 lines. Linux is around 30 million. The difference is not
conceptual --- it is drivers, architectures, filesystems, protocols, decades of edge cases, and
performance work.

The concepts you have now are the same ones. A page table in Linux is a page table here. A context
switch there saves the same registers. A system call crosses the same boundary in the same way. When
you read that a CPU vulnerability breaks the isolation between user space and kernel space, you now
know exactly what that sentence means and exactly which bit is involved.

That is what this was for.

\begin{summary}
You started not knowing what a register was. You now understand a complete operating system, you
can read a processor trace and find a memory corruption bug from it, and you have a list of real
improvements you could make to a real codebase. The appendices that follow are reference material
for when you do.
\end{summary}
